% page 16 of the facsimile (image p-015 of the scan)
% block 154.9 x 229.0 mm ; left margin 8.0 mm
\pgc{-12.700mm}{0.000mm}{
\l{\cel{5}8.}
\l{}
\l{}
\l{\sou{adversa}. Kontrea. (\sou{En lukto}) La persono qua opozesas ad}
\l{\cel{3}altra. - EFIS.}
\l{}
\l{\sou{advokato}. La homo di qua la profesiono konsistas ye pledar}
\l{\cel{3}por la personi qui procesas koram tribunalo. - DEFIRS.}
\l{}
\l{\sou{adyunto}. Qua juntesas ad ulu kom lua helpanto. - DEFIRS.}
\l{}
\l{\sou{aero}. Fluido gasa, diafana, ek oxo, azoto e kelka gasi skar-}
\l{\cel{3}sa, difuzita en omna regioni di Tero, utila por respiro}
\l{\cel{3}e kombusto. - DEFIS.}
\l{}
\l{\sou{aerenkimo}. (\sou{biol.}) Tisuo ek lakuni respirala, quan onu ren-}
\l{\cel{3}kontras en la organi (precipue la radiki) qui kreskas en}
\l{\cel{3}slamo.}
\l{}
\l{\sou{aerobia}. (\sou{biol}.)\cel{2}Qua bezonas la oxo di aero por respirar.}
\l{\cel{3}- DEFIS.}
\l{}
\l{\sou{aerolito}. (\sou{astron.})\cel{2}Maso minerala qua falas adsur la surfa-}
\l{\cel{3}co di Tero, de la regioni superega di atmosfero. - DEFIRS.}
\l{}
\l{\sou{aerometro}. Instrumento per qua onu povas mezurar la dens-}
\l{\cel{3}eso-grado di aero, di gasi. - DEFIRS.}
\l{}
\l{\sou{aerometrio.}\cel{2}La parto di fiziko qua studias la denseso-grado}
\l{\cel{3}di aero, di gasi. - DEFIRS.}
\l{}
\l{\sou{aeroplano}. (\sou{aviaco}) Mashino qua povas sustenar su super-sule,}
\l{\cel{3}en atmosfero, sen esar plu lejera kam ica, efike da la}
\l{\cel{3}preso da vento sur la surfaci inklinita di lua quaza ali.}
\l{\cel{3}- DEFIRS.}
\l{}
\l{\sou{aeroskopo}. Instrumento per qua onu kolektas (por studiar li)}
\l{\cel{3}la polvuni mikroskopala qui suspendesas en atmosfero. -}
\l{\cel{3}DEFIRS.}
\l{}
\l{\sou{aerostato}.\cel{2}Balonego di qua la susteneso en atmosfero obten-}
\l{\cel{3}esas per la uzo di gaso plu lejera kam aero. - DEFIRS.}
\l{}
\l{\sou{afabla}.- Qua aceptas sua kunhomi benigne, aminde. - EF IS.}
\l{}
\l{\sou{afazio}. (\sou{patol}.) Privaceso di la parol-fakultato, konjekt-}
\l{\cel{3}eble pro lezuro cerebrala. - DEFIS.}
\l{}
\l{\sou{afecionar}. (\sou{trans}.) Unionesar ad (ulu, ulo) per amo o prizo}
\l{\cel{3}duranta. - EFIS.}
\l{}
\l{\sou{afektar}\cel{2}(\sou{trans}.) I. (\sou{medicino}) (\sou{pri morbo}) Efikar a la homo-}
\l{\cel{3}korpo. - II. (\sou{psikol.}) Efikar a la psiko. - DEFIRS.}
}
